\documentclass[11pt]{article}

\usepackage[margin=1in]{geometry}
\usepackage{amsmath,amssymb}
\usepackage{booktabs}
\usepackage[hidelinks]{hyperref}

\title{Symmetry of Initial-$N_J$ Cutoffs for Squeezing}
\author{}
\date{}

\begin{document}
\maketitle

\section{Conclusion}

Reflection symmetry of the $N_J$ distribution about $N/2$ is not required for
one-axis twisting (OAT) to generate squeezing.  It is, however, part of the
canonical initial coherent spin state used to derive the usual OAT formulas
and the unsqueezed baseline $\xi^2(0)=1$.  A substantial asymmetric truncation
therefore changes the initial state and its squeezing dynamics, even though it
does not remove the OAT mechanism.

For enforcing $\lvert\Omega/\Omega_c\rvert<r_{\max}$, retain the valid large
$N_J$ sectors and remove only the invalid low-$N_J$ sectors.  A mirrored
high-$N_J$ cutoff has no physical justification: since
\begin{equation}
  \frac{\lvert\Omega\rvert}{\Omega_c(N_J)}
  = \frac{2\lvert\Omega\rvert}{\Gamma N_J},
\end{equation}
larger $N_J$ sectors lie farther inside the spin-polarized phase.  The
one-sided cutoff is safe as an approximation only when its discarded
probability is negligible.  In an inhomogeneous model, the same argument
applies to the corresponding positive coupling-weighted active population.

\section{Theory Basis}

The paper starts from the equal single-particle superposition, whose $N_J$
probabilities are binomial and symmetric about $N/2$.  Sector-dependent Berry
phases then shear this coherent-state distribution.  Writing fluctuations
about a general local center $\bar N_J$,
\begin{equation}
  \phi_B(N_J)
  = \phi_0 + \phi_1\Delta N_J
    + \frac{1}{2}\phi_2(\Delta N_J)^2 + \cdots,
  \qquad \Delta N_J=N_J-\bar N_J,
\end{equation}
shows that the constant is a global phase, the linear term is a rotation, and
the quadratic curvature produces squeezing.  Exact reflection symmetry does
not enter this mechanism.  What matters is a localized distribution with
enough variance and coherence between nearby sectors for the quadratic
description to apply.

The standard analytic result expands about $N_J=N/2$, or equivalently
$S_z=0$, and the Holstein--Primakoff calculation takes equal $J$- and
$\uparrow$-state fractions.  A material one-sided cut produces nonzero mean
$S_z$, skewness, and a non-coherent initial covariance.  Its effects vanish
with the discarded conditional probability $q$; for an originally centered
distribution, the induced mean is
\begin{equation}
  \langle S_z\rangle_{\mathrm{kept}}
  = -\frac{q}{1-q}
    \left(\frac{N}{2}
    - \mathbb E[N_J\mid\text{discarded}]\right).
\end{equation}

The symmetric $N/2\pm dN$ window used in the paper is therefore a convenient
numerical truncation of a centered binomial distribution, not a dynamical OAT
requirement.  Squeezing depends on local, low-order sector coherences rather
than on pairing every $N_J$ sector with its mirror $N-N_J$.

\section{Current Constrained Scan}

For the current $N_i=(50,50)$ scan,
$\Omega_{\max}=0.1N\Gamma=10\Gamma$.  With the present homogeneous-form code
check, the strict limit $\lvert\Omega/\Omega_c\rvert<1$ requires $N_J>20$.
The inhomogeneous mean-field equation instead suggests the sector-wise
weighted extension
\begin{equation}
  N_{J,\omega}=0.7N_{J1}+1.3N_{J2}>20.
\end{equation}
This extension should be validated separately before replacing the current
check; the project theory note establishes the homogeneous form explicitly
only at the central weighted sector.  The table uses exact product-binomial
weights conditioned on each line $N_{J1}=A N_{J2}$.  Here $q_{\rm tot}$ is
removed by the current $N_J\leq20$ rule and $q_\omega$ by the weighted rule.

\begin{center}
\begin{tabular}{@{}rrrrr@{}}
\toprule
$A$ & $\mathbb E[N_J]$ & $\sigma(N_J)$
& $q_{\rm tot}$ & $q_\omega$ \\
\midrule
$1$   & $50.000$ & $5.025$ & $1.11\times10^{-7}\%$ & $1.11\times10^{-7}\%$ \\
$1.5$ & $48.087$ & $4.842$ & $1.12\times10^{-6}\%$ & $1.12\times10^{-6}\%$ \\
$2$   & $45.178$ & $4.597$ & $4.76\times10^{-7}\%$ & $2.33\times10^{-5}\%$ \\
$3$   & $40.650$ & $4.278$ & $3.63\times10^{-4}\%$ & $3.63\times10^{-4}\%$ \\
$4$   & $37.729$ & $4.104$ & $4.93\times10^{-3}\%$ & $4.93\times10^{-3}\%$ \\
\bottomrule
\end{tabular}
\end{center}

Thus, either low-only validity estimate is negligible for all five cases;
even the largest change removes only about $0.005\%$ of the corresponding
constraint-conditioned state.  A symmetric mirror cutoff would make no
meaningful numerical difference here, but it would discard valid states for
no theoretical benefit.

The larger conceptual change already comes from the line constraints.  After
renormalization they define projected, generally non-coherent initial states;
for example, the conditional mean moves from $50$ at $A=1$ to about $37.7$ at
$A=4$.  Differences in squeezing across $A$ therefore include differences in
the initial state, not only differences in subsequent dynamics.  A generalized
squeezing metric remains computable, but canonical coherent-state OAT formulas
need not apply quantitatively.

\section{Recommendation}

Use a lower cutoff derived from the requested maximum
$\lvert\Omega/\Omega_c\rvert$, without a mirrored upper cutoff.  Report, for
each constrained initial state, the discarded conditional probability,
$\mathbb E[N_J]$, $\sigma(N_J)$, skewness, and $\xi^2(0)$.  If the cutoff
removes appreciable probability or approaches the central support, reduce
$\lvert\Omega\rvert$ instead of repairing the distribution with a symmetric
cut: neither truncation would then faithfully represent the intended initial
state.

\section{Local Sources}

The conclusions use:
\begingroup
\raggedright
\begin{itemize}
  \item the initial coherent state, strong-symmetry sector evolution, and OAT
    expansion in \path{docs/theory/main.tex};
  \item the Berry-phase expansion and symmetric numerical benchmark window in
    \path{docs/theory/end_matter.tex};
  \item the equatorial fluctuation expansion in
    \path{docs/theory/appendix_holstein_primakoff.tex};
  \item the nearby-sector coherence requirement in
    \path{docs/theory/appendix_coherence_preservation.tex}.
  \item the group-weighted critical-drive relation in
    \path{docs/theory/notes_inhomogeneous_coupling.tex}.
\end{itemize}
\endgroup

\end{document}
